\section{Axiom Verification}
\label{sec:axioms}

The foundational paper's structural results (self-balancing, scorpion
detection, cooperative expansion) hold for any measure satisfying axioms
M1--M6. This section proves that $\volL$ satisfies all six.

\begin{proposition}[Poset Measure Satisfies M1--M6]
\label{prop:axioms}
The poset measure $\volL$ (Definition~\ref{def:poset_measure}) satisfies
    the following axioms for all measurable $G, A, B \subseteq \C$:
\begin{enumerate}
\item[\textbf{(M1)}] $\volL(G) \geq 0$.
\item[\textbf{(M2)}] $\volL(\emptyset) = 0$.
\item[\textbf{(M3)}] $A \subseteq B \implies \volL(A) \leq \volL(B)$.
\item[\textbf{(M4)}] If $A$ and $B$ are \emph{poset-disjoint}
    (Definition~\ref{def:poset_disjoint}), then
    $\volL(A \cup B) = \volL(A) + \volL(B)$.
\item[\textbf{(M5)}] For every benchmarked $d \in \C$ with $\smax(d) \geq 1$,
    $\volL(\{d\}) > 0$.
\item[\textbf{(M6)}] When agents' individual capability sets occupy
    incomparable regions of the poset, $\volL(G) \geq \sum_k
    \volL(\Gind{k})$, with strict inequality when cooperative capabilities
    exist.
\end{enumerate}
\end{proposition}

Before proving the proposition, we define the disjointness condition that M4
requires on the poset.

\begin{definition}[Poset-Disjoint]
\label{def:poset_disjoint}
Two sets $A, B \subseteq \C$ are \emph{poset-disjoint} if no element of $A$
    subsumes or is subsumed by any element of $B$: for all $a \in A$ and $b
    \in B$, neither $a \leq b$ nor $b \leq a$. Equivalently, $A$ and $B$
    occupy incomparable regions of the poset, and their downward closures
    intersect only at $\bot$: $\dc A \cap \dc B = \{\bot\}$.
\end{definition}

This is stronger than set-disjointness ($A \cap B = \emptyset$). Two
capabilities can be set-disjoint (they are different nodes) but not
poset-disjoint (one subsumes the other). Poset-disjointness ensures that
the capabilities in $A$ and $B$ are structurally independent in the
subsumption hierarchy.

\begin{proof}[Proof of Proposition~\ref{prop:axioms}]
\emph{(M1) Non-negativity.} By the design constraint on $w$
(Section~\ref{sec:poset}), $\iw(d) \geq 0$ for all $d \in \C$. Since
$\volL(G) = \sum_{d \in \dc G} \iw(d)$ is a sum of non-negative terms,
$\volL(G) \geq 0$.

\emph{(M2) Null empty set.} $\dc \emptyset = \emptyset$ (the downward closure
of the empty set is empty, since no element of $\C$ is below any element of
$\emptyset$). The sum $\volL(\emptyset) = \sum_{d \in \emptyset} \iw(d) = 0$
(empty sum).

\emph{(M3) Monotonicity.} If $A \subseteq B$, then $\dc A \subseteq \dc B$
(the downward closure of a subset is a subset of the downward closure). Since
every term in the sum $\volL(A) = \sum_{d \in \dc A} \iw(d)$ also appears in
$\volL(B) = \sum_{d \in \dc B} \iw(d)$, and all terms are non-negative by
(M1), $\volL(A) \leq \volL(B)$.

\emph{(M4) Finite additivity under poset-disjointness.} If $A$ and $B$ are
poset-disjoint, then $\dc A \cap \dc B = \{\bot\}$. The downward closure of
the union satisfies $\dc(A \cup B) = \dc A \cup \dc B$. Since the only shared
element is $\bot$ (with $\iw(\bot) = 0$):
\[
\volL(A \cup B) = \sum_{d \in \dc A \cup \dc B} \iw(d)
  = \sum_{d \in \dc A} \iw(d) + \sum_{d \in \dc B} \iw(d)
  - \iw(\bot)
  = \volL(A) + \volL(B).
\]

\emph{(M5) Non-triviality.} For any atom $d$ of the poset (a minimal
non-$\bot$ element), $\iw(d) = w(d) - w(\bot) = w(d) > 0$ since $\smax(d)
\geq 1$ implies $w(d) = \log_2(1 + \smax(d)) \geq \log_2(2) = 1$. For
non-atoms, $\volL(\{d\}) = \sum_{d' \in \dc\{d\}} \iw(d')$. By M1
(all $\iw(d') \geq 0$, verified by Lemma~\ref{lem:nonneg_iw}), no term in
this sum is negative, and the sum includes at least one atom with
$\iw > 0$, so $\volL(\{d\}) > 0$.

\emph{(M6) Superadditivity under independence.} When agents' individual
capability sets $\Gind{1}, \ldots, \Gind{n}$ occupy pairwise incomparable
regions of the poset (no agent's capability subsumes another agent's
capability), the sets are pairwise poset-disjoint. By (M4) applied
iteratively, $\volL(\bigcup_k \Gind{k}) = \sum_k \volL(\Gind{k})$. Adding
cooperative capabilities $\Gcoop$ to form $G = \bigcup_k \Gind{k} \cup
\Gcoop$: since $\Gcoop$ contains capabilities not in any individual set,
$\dc G \supsetneq \dc(\bigcup_k \Gind{k})$, and by (M3),
$\volL(G) \geq \volL(\bigcup_k \Gind{k}) = \sum_k \volL(\Gind{k})$.

If $\Gcoop \neq \emptyset$ and contains at least one capability $d_c$ not
subsumed by any element of $\bigcup_k \Gind{k}$, then $\dc G$ contains
$d_c \notin \dc(\bigcup_k \Gind{k})$, and $\iw(d_c) > 0$ by (M5), giving
strict inequality: $\volL(G) > \sum_k \volL(\Gind{k})$.
\end{proof}

\begin{theorem}[Self-Balancing on Finite Posets]
\label{thm:self_balancing_poset}
The self-balancing property (Proposition~1 of \citet{lasser2026gfm}) holds for
    $\volL$ on finite weighted posets. Specifically: destructive
    permissiveness, coercive restriction, and self-imposed rigidity are all
    $\volL$-contracting under the conditions of Lemmas~1--3 of the foundational
    paper.
\end{theorem}

\begin{proof}
By Proposition~\ref{prop:axioms}, $\volL$ satisfies M1--M6. The foundational
paper's Lemmas~1--3 and Proposition~1 are derived from M1--M6 alone (see
Appendix~E of \citet{lasser2026gfm}). Therefore all structural results hold
with $\vol$ replaced by $\volL$ throughout.
\end{proof}
